\label{non-abelian-cech}

\rednote{TODO Are $X$ and $Y$ really an arbitrary type or did I use they were sets at some point?}

In this section we work in plain HoTT and study non-abelian torsors over an arbitrary type $X$. Our goal is to describe non-abelian torsors over $X$ trivialised by a given cover in terms of cocycles for this cover. Recall that given a group $G$, a $G$-pseudotorsor is a type with a transitive and free right action of $G$ on $X$, and a $G$-torsor is an inhabited $G$-pseudotorsor.

\begin{definition}
Assume given $G$ a group and $\alpha$ a $G$-torsor with $x,y:\alpha$. Then $x^{-1}\cdot y$ is defined as the only $g:G$ such that $x\cdot g = y$. 
\end{definition}

Note that $(x^{-1}\cdot y)\cdot (y^{-1}\cdot z) = x^{-1}\cdot z$. Now we define non-abelian \v{C}ech cohomology.

\begin{definition}\label{dfnCocycle}
Assume given $X$ a type with $Y(x)$ an inhabited type and $G(x)$ a group both depending on $x:X$.
\begin{itemize}
\item A cocycle is a $\delta:\Pi_{x:X}G(x)^{Y(x)^2}$ such that for all $x:X$ and $y_0,y_1,y_2:Y(x)$ we have $\delta_x(y_0,y_1)\cdot \delta_x(y_1,y_2) = \delta_x(y_0,y_2)$.
\item Two cocycles $\delta$ and $\delta'$ are called cohomologous if there exists $\gamma:\Pi_{x:X}G(x)^{Y(x)}$ such that for all $x:X$ and $y_0,y_1:Y(x)$ we have $\gamma_x(y_0)\cdot\delta_x(y_0,y_1) = \delta'_x(y_0,y_1)\cdot\gamma_x(y_1)$.
\end{itemize}
\end{definition}

Being cohomologous in an equivalence relation. 

\begin{definition}\label{dfnCechCohomology}
Assume given $X$ a type with $Y(x)$ an inhabited type and $G(x)$ a group both depending on $x:X$. We define the non-abelian cover cohomology $\check{H}^1(X,Y,G)$ as the quotient of cocycles by being cohomologous.
\end{definition}

We prove a first easy lemma.

\begin{lemma}\label{first-cech-has-section-vanish}
Assume given $X$ a type with $Y(x)$ an inhabited type and $G(x)$ a group both depending on $x:X$. If we merely have $\Pi_{x:X}Y(x)$ then $\check{H}^1(X,Y,G)$ is a singleton.
\end{lemma}

\begin{proof}
  Assume given $t: \Pi_{x:X} Y(x)$ and $\delta:\Pi_{x:X}G(x)^{Y(x)^2}$ a cocycle, we need to show it is cohomologous to $1$. It is enough to consider $\gamma:\Pi_{x:X}G(x)^{Y(x)}$ defined by $\gamma_x(y) = \delta_x(t_x,y)$.
\end{proof}

Our main goal is now to show that the cover cohomology $\check{H}^1(X,Y,G)$ corresponds to $G$-torsors over $X$ that are trivialised by $Y$. We state this precisely using the following auxiliary definition.

\begin{definition}
Assume given $X$ a type with $Y(x)$ an inhabited type and $G(x)$ a group both depending on $x:X$. We define $H^1(X,Y,G)$ as $\{\alpha:H^1(X,G)\ |\ \pi_X^*(\alpha) = 0 \}$ where $\pi_X$ is the projection $\Sigma_{x:X}Y(x)\to X$.
\end{definition}

Our goal is then to show $H^1(X,Y,G) = \check{H}^1(X,Y,G)$. First we define a map from $H^1(X,Y,G)$ to $\check{H}^1(X,Y,G)$.

\begin{lemma}\label{phi-well-defined}
Assume given $X$ a type with $Y(x)$ an inhabited type and $G(x)$ a group both depending on $x:X$. Given $\alpha:\Pi_{x:X}BG(x)$ and $\beta:\Pi_{x:X}\alpha(x)^{Y(x)}$, we define $\phi(\alpha,\beta):\Pi_{x:X}G(x)^{Y(x)^2}$ by $\phi(\alpha,\beta)_x(y_0,y_1) = \beta_x(y_0)^{-1}\cdot\beta_x(y_1)$. Then we have that:

\begin{enumerate}[(i)]
\item $\phi(\alpha,\beta)$ is a cocycle.
\item Given another $\beta':\Pi_{x:X}\alpha(x)^{Y(x)}$, we have that $\phi(\alpha,\beta)$ and $\phi(\alpha,\beta')$ are cohomologous.
\item We have an induced map $\Phi : H^1(X,Y,G)\to \check{H}^1(X,Y,G)$ defined by $\Phi(\alpha) = \phi(\alpha,\beta)$ for any $\beta:\Pi_{x:X}\alpha(x)^{Y(x)}$.
\end{enumerate}
\end{lemma}

\begin{proof}
As follows:

\begin{enumerate}[(i)]
\item Clear.
\item To show that $\phi(\alpha,\beta)$ and $\phi(\alpha,\beta')$ are cohomologous, we consider $\gamma = \lambda x,y.\, \beta'_x(y)^{-1}\cdot\beta_x(y)$. Given $x:X$ and $y_0,y_1:Y(x)$ we indeed have:
\[\phi(\alpha,\beta')_x(y_0,y_1)\cdot\gamma_x(y_1) = (\beta'_x(y_0)^{-1}\cdot \beta'_x(y_1))\cdot (\beta'_x(y_1)^{-1}\cdot\beta_x(y_1) = \beta'_x(y_0)^{-1}\cdot\beta_x(y_1)\]
\[\gamma_x(y_0)\cdot\phi(\alpha,\beta)_x(y_0,y_1) = (\beta'_x(y_0)^{-1}\cdot \beta_x(y_0))\cdot(\beta_x(y_0)^{-1}\cdot\beta_x(y_1)) = \beta'_x(y_0)^{-1}\cdot\beta_x(y_1)\]
\item Since $\alpha$ is trivialised by $\pi_X$, we know that there there exists such a $\beta$. By (ii) the class of $\psi(\alpha,\beta)$ in $\check{H}^1(X,Y,G)$ does not depend on the choice of such a $\beta$.
\end{enumerate}
\end{proof}

Now we define a map from $\check{H}^1(X,Y,G)$ to $H^1(X,Y,G)$.

\begin{lemma}\label{psi-well-defined}
Assume given $Y$ an inhabited type and $G$ a group, as well as $\delta:G^{Y^2}$ a cocycle. %such that for all $y_0,y_1,y_2:Y$ we have $\delta(y_0,y_1)\cdot\delta(y_1,y_2) = \delta(y_0,y_2)$. 
Then we define $\psi(\delta) = \{\gamma:G^Y\ |\ \forall y_0,y_1.\, \gamma(y_0)\cdot\gamma(y_1)^{-1} = \delta(y_0,y_1)\}$. Then we have that:

\begin{enumerate}[(i)] 
\item We have a $H:\psi(\delta)^Y$ defined by $H(y) = \lambda y'.\, \delta(y',y)$.
\item The type $\psi(\delta)$ is a $G$-torsor.
\item Given $\delta$ and $\delta'$ cocycles with $\iota:G^Y$ such that for all $y_0,y_1:Y$ we have $\iota(y_0)\cdot\delta(y_0,y_1) = \delta(y_0,y_1)\cdot\iota(y_1)$, we have $\psi(\delta)=\psi(\delta')$.
\item We have an induced map $\Psi : \check{H}^1(X,Y,G)\to H^1(X,Y,G)$ defined by $\Psi(\delta) = \lambda x.\, \psi(\delta_x)$.
\end{enumerate}
\end{lemma}

\begin{proof}
We prove this as follows:

\begin{enumerate}[(i)]
\item We just need to check that for all $\bar{y},y_0,y_1$ we indeed have $\delta(y_0,\bar{y})\cdot\delta(y_1,\bar{y})^{-1} = \delta(y_0,y_1)$, but this follows from $\delta$ being a cocycle.

\item We define the action of $g:G$ on $\gamma:\psi(\delta)$ by $\gamma\cdot g = \lambda y.\, \gamma(y)\cdot g$.

It is clear that the action is well-defined and free, now let us check it is transitive. Assume $\gamma,\gamma':\psi(\delta)$, then we can check that $\gamma(y)^{-1}\cdot\gamma'(y)$ does not depend on $y$, indeed $\gamma(y_0)\cdot\gamma(y_1)^{-1} = \delta_x(y_0,y_1) = \gamma'(y_0)\cdot\gamma'(y_1)^{-1}$. Since $Y$ is inhabited this gives us an element $g:G$ such that $\gamma(y)\cdot g = \gamma'(y)$ for all $y:Y$, which is precisely what we want.

The fact that $\propTrunc{\psi(\delta)}$ follows directly from $Y$ inhabited and (i).

\item The map in $\psi(\delta)\to \psi(\delta')$ sending $\gamma$ to $\lambda y.\, \iota(y)\cdot\gamma(y)$ is an equivalence of torsors.

\item Given (ii) and (iii) it is clear that we get an element in $H^1(X,G)$. From $(i)$ we see that we have $\Pi_{x:X}\psi(\delta_x)^{Y(x)}$, so that $\lambda x.\, \psi(\delta_x)$ is indeed trivialised by $\pi$.
\end{enumerate}
\end{proof}

Now we prove the two maps are inverse to each other.

\begin{proposition}\label{non-abelian-chech-is-sheaf-generic}
Assume given $X$ a type with $Y(x)$ an inhabited type and $G(x)$ a group both depending on $x:X$. Then $H^1(X,Y,G) = \check{H}^1(X,Y,G)$.
\end{proposition}

\begin{proof}
First we check that $\Psi\circ\Phi$ is the identity. It is enough to assume $\alpha:\Pi_{x:X}BG(x)$ with $\beta:\Pi_{x:X}\alpha(x)^{Y(x)}$ and prove that for all $x:X$ we have 
\[\alpha(x) = \psi(\phi(\alpha,\beta)_x)\]
where
\[\psi(\phi_x(\alpha,\beta)) = \{\gamma:G(x)^{Y(x)}\ |\ \forall y_0,y_1.\, \gamma(y_0)\cdot\gamma(y_1)^{-1} = \beta_x(y_0)^{-1}\cdot\beta_x(y_1)\}\]

We define $f:\alpha(x)\to \psi(\phi(\alpha,\beta)_x)$ by $f(e) = \lambda y.\, \beta_x(y)^{-1}\cdot e$. This is well defined as for all $y_0,y_1:Y(x)$ we have that:
\[f(p,y_0)\cdot f(p,y_1)^{-1} = \beta_x(y_0)^{-1}\cdot e \cdot (\beta_x(y_1)^{-1}\cdot e)^{-1} = \beta_x(y_0)^{-1}\cdot\beta_x(y_1)\]
Morever this map clearly commutes with the group action so it is enough to show it is an equivalence of type to conclude.

In the other direction, assume given $\gamma:\psi(\phi(\alpha,\beta)_x)$. Then we have that $\beta_x(y)\cdot\gamma(y):\alpha(x)$ does not depend on the choice of $y:Y(x)$ by the hypothesis on $\gamma$, and since $Y(x)$ is inhabited this defines an element $g(\gamma):\alpha(x)$.

Next we check that $g(f(e)) = e$, since we want to prove a proposition we can assume $y:Y(x)$, then we need to prove that $\beta_x(y)\cdot (\beta_x(y)^{-1}\cdot e) = e$ which is clear. Finally to check $f(g(\gamma)) = \gamma$ we need to prove that $\beta_x(y)^{-1}\cdot \beta_x(y) \cdot \gamma(y) = \gamma(y)$ which is also obvious (we used the fact that $g(\gamma)$ can be computed using any $y:Y(x)$).

Now we check that $\Phi\circ \Psi$ is the identity. It is enough to prove that given a cocycle $\delta$, for all $x:X$ with $y_0,y_1:Y(x)$ we have that $H_x(y_0)^{-1}\cdot H_x(y_1) = \delta_x(y_0,y_1)$ where $H_x : Y(x)\to \psi(\delta_x)$ is defined in \Cref{psi-well-defined}, point (ii). But recall that $H_x(\bar{y}) = \lambda y.\delta_x(y,\bar{y})$ so we need to check that $(\lambda y.\delta_x(y,y_0))\cdot\delta_x(y_0,y_1) = \lambda y.\delta_x(y,y_1)$ which is clear as the action is defined pointwise.
\end{proof}

\begin{corollary}\label{non-abelian-chech-is-sheaf-generic-ayclic-cover}
Assume given $X$ a type with $Y(x)$ an inhabited type and $G(x)$ a group both depending on $x:X$. If $H^1(\Sigma_{x:X}Y(x), G\circ\pi_X) = 0$, then $H^1(X,G) = \check{H}^1(X,Y,G)$.
\end{corollary}
